\chapter{Conteúdo Teórico}
\label{cap:teorico}

Este capítulo reúne os conceitos necessários para acompanhar as escolhas
metodológicas e a leitura dos resultados. A Seção~\ref{sec:T-complexidade}
apresenta a noção de complexidade estatística e a formulação LMC. A
Seção~\ref{sec:T-sampen} define a entropia amostral e sua extensão
bidimensional. A Seção~\ref{sec:T-camada} justifica a escolha da camada densa
como objeto de medida e formaliza a leitura de seus pesos como sequência. A
Seção~\ref{sec:T-overfitting} fixa a definição operacional de
\emph{overfitting} adotada.

\section{Complexidade Estatística}
\label{sec:T-complexidade}

\subsection{Entropia de Shannon e a insuficiência da desordem}

Dada uma variável discreta com distribuição de probabilidade
$P = \{p_1, \dots, p_N\}$, a entropia de Shannon \cite{shannon1948} é

\begin{equation}
    H(P) = -\sum_{i=1}^{N} p_i \ln p_i .
    \label{eq:shannon}
\end{equation}

A entropia mede desordem: é nula quando toda a probabilidade concentra-se em um
único estado e máxima, igual a $\ln N$, quando a distribuição é uniforme. Por
isso, isoladamente, ela não serve como medida de complexidade. Um cristal
perfeito tem entropia mínima e um gás ideal tem entropia máxima, mas nenhum dos
dois é estruturalmente complexo: o primeiro é totalmente previsível, o segundo
totalmente aleatório. A complexidade que interessa está entre os dois extremos.

\subsection{Complexidade LMC}

\citeonline{lmc1995} resolvem esse problema introduzindo um segundo termo, o
desequilíbrio, que mede a distância quadrática da distribuição observada em
relação à distribuição uniforme $P_e = \{1/N, \dots, 1/N\}$:

\begin{equation}
    D(P) = \sum_{i=1}^{N} \left( p_i - \frac{1}{N} \right)^{2}.
    \label{eq:desequilibrio}
\end{equation}

A complexidade LMC é o produto dos dois termos, com a entropia normalizada pelo
seu valor máximo:

\begin{equation}
    C_{\mathrm{LMC}}(P) = \frac{H(P)}{\ln N} \cdot D(P).
    \label{eq:lmc}
\end{equation}

O produto na Equação~\ref{eq:lmc} anula-se nos dois extremos e é máximo em
distribuições intermediárias. No cristal, $H \to 0$ e o primeiro fator zera; no
gás ideal, $P \to P_e$ e o desequilíbrio zera. Uma discussão das propriedades e
das limitações dessa formulação, incluindo variantes baseadas na divergência de
Jensen-Shannon, encontra-se em \citeonline{lamberti2004} e
\citeonline{rosso2010}.

Duas propriedades são relevantes para este trabalho. A primeira é que
$C_{\mathrm{LMC}}$ depende apenas do histograma dos valores, sendo portanto
\emph{invariante a qualquer permutação} da sequência analisada. A segunda,
consequência da primeira, é que a LMC não pode captar nenhuma informação
posicional: se a memorização se manifestar na organização espacial dos pesos e
não na distribuição de seus valores, a LMC será cega a ela.

\section{Entropia Amostral}
\label{sec:T-sampen}

\subsection{Formulação unidimensional}

A entropia amostral \cite{richman2000} quantifica a irregularidade de uma
sequência medindo com que frequência padrões semelhantes permanecem semelhantes
quando estendidos por uma amostra. Ela corrige o viés de autocomparação
presente na entropia aproximada de \citeonline{pincus1991}.

Seja a série $\{u(1), \dots, u(L)\}$ e os vetores de janela
$\mathbf{x}_m(i) = \{u(i), \dots, u(i+m-1)\}$. Sejam $B$ o número de pares de
janelas de comprimento $m$ cuja distância de Chebyshev é inferior à tolerância
$r$, e $A$ o número correspondente para janelas de comprimento $m+1$. Define-se

\begin{equation}
    \mathrm{SampEn}(m, r, L) = -\ln \frac{A}{B}.
    \label{eq:sampen}
\end{equation}

O parâmetro $m$ é o comprimento do padrão comparado e $r$ é a tolerância de
similaridade, usualmente expressa como fração do desvio-padrão da série. Valores
elevados de $\mathrm{SampEn}$ indicam baixa previsibilidade. A escolha de $m$ e
$r$ não é neutra e afeta a estabilidade da estimativa, especialmente em séries
curtas \cite{yentes2013}; por isso este trabalho reporta a sensibilidade dos
resultados a esses parâmetros.

Diferentemente da LMC, a entropia amostral \emph{não} é invariante à ordem, pois
compara janelas consecutivas. Essa é precisamente a propriedade que a torna
complementar à LMC neste trabalho.

\subsection{Extensão bidimensional}

A entropia amostral bidimensional \cite{silva2016} generaliza a
Equação~\ref{eq:sampen} substituindo as janelas unidimensionais por
sub-matrizes $m \times m$ deslizantes sobre uma matriz, com a mesma lógica de
contagem de pares similares. Preserva-se assim a vizinhança nas duas direções,
em vez de destruí-la pelo achatamento. Sua formulação e validação em textura de
imagens são discutidas em \citeonline{silva2014embc} e \citeonline{silva2012qsampen}.

Para este trabalho a SampEn2D é particularmente pertinente, pois a matriz de
pesos de uma camada densa é um objeto naturalmente bidimensional --- cada linha
corresponde a um neurônio de saída e cada coluna a uma característica de
entrada --- e reduzi-la a um vetor descarta uma das duas estruturas.

\subsection{Entropia multiescala}

A entropia multiescala \cite{costa2002, costa2005} aplica a entropia amostral a
versões da série granuladas em escalas crescentes, obtidas pela média de blocos
consecutivos não sobrepostos. O perfil resultante distingue sinais cuja
irregularidade se mantém através das escalas daqueles em que ela se concentra
na escala fina. Está implementada neste trabalho e prevista como extensão.

\section{A Camada Densa como Objeto de Medida}
\label{sec:T-camada}

\subsection{Por que a camada densa}

Em uma rede convolucional de classificação, as camadas convolucionais produzem
uma representação da imagem e a camada densa final converte essa representação
em pontuações por classe. É nela que a decisão é efetivamente tomada, e é nela
que a associação entre representação e rótulo precisa ser armazenada. Se a rede
memoriza pares imagem--rótulo inconsistentes, é razoável esperar que essa
associação se inscreva de forma mais direta nos pesos dessa camada. A escolha
tem ainda a vantagem prática de manter o custo do monitoramento baixo: na
ResNet-18 aplicada a seis classes, a matriz de pesos tem forma $6 \times 512$,
totalizando \num{3072} parâmetros.

\subsection{O caminho e a ordem de achatamento}

Considere-se uma camada densa com $M$ entradas $x_1, \dots, x_M$ e $N$ neurônios
$n_1, \dots, n_N$. A matriz de pesos é $W \in \mathbb{R}^{N \times M}$, e o
elemento $w_{ji}$ é o peso do \emph{caminho} $x_i \rightarrow n_j$. Essa é
exatamente a convenção do PyTorch, em que o atributo de pesos de uma camada
linear tem forma $[N, M]$.

Como a entropia amostral unidimensional depende da ordem, é necessário fixar o
percurso pelo qual a matriz é lida como sequência. Duas ordens são naturais.
Na ordem \emph{por neurônio}, percorrem-se todos os pesos que chegam a cada
neurônio antes de passar ao seguinte:

\begin{equation}
    v_{n} = (w_{11}, w_{12}, \dots, w_{1M},\; w_{21}, \dots, w_{2M},\; \dots,\; w_{NM}).
    \label{eq:ordem-n}
\end{equation}

Na ordem \emph{por entrada}, percorrem-se todos os pesos que partem de cada
entrada:

\begin{equation}
    v_{x} = (w_{11}, w_{21}, \dots, w_{N1},\; w_{12}, \dots, w_{N2},\; \dots).
    \label{eq:ordem-x}
\end{equation}

As duas leituras não são equivalentes e medem coisas distintas: a
Equação~\ref{eq:ordem-n} avalia a regularidade \emph{dentro do campo receptivo
de cada neurônio}, enquanto a Equação~\ref{eq:ordem-x} avalia a regularidade
\emph{da distribuição de uma mesma entrada entre os neurônios}. Adota-se a
ordem por neurônio como convenção padrão, e a comparação entre as duas é
reportada como ablação. Para a complexidade LMC a distinção é irrelevante, pela
invariância discutida na Seção~\ref{sec:T-complexidade}.

Por coerência com a noção de caminho, o vetor de viés é excluído da análise: o
viés é um deslocamento aditivo do neurônio, não a ligação entre uma entrada e
um neurônio. Quando há mais de uma camada densa, cada uma é analisada
separadamente, sem concatenação.

\section{Definição Operacional de \emph{Overfitting}}
\label{sec:T-overfitting}

Para que a antecedência de um indicador seja mensurável, é preciso uma
definição de \emph{overfitting} que produza um número de época, e não apenas
uma descrição qualitativa. Adota-se a seguinte construção.

Seja $\mathrm{val\_loss}(e)$ a perda no conjunto de validação na época $e$. A
melhor época é $e^{*} = \arg\min_{e} \mathrm{val\_loss}(e)$. O início do
\emph{overfitting} é a primeira época $e^{of} > e^{*}$ tal que
$\mathrm{val\_loss}(e) > \mathrm{val\_loss}(e^{*}) + \delta$ de forma sustentada
por um número pré-fixado de épocas consecutivas. Toda análise dos indicadores
restringe-se às épocas $e \le e^{of}$; as posteriores permanecem registradas e
são exibidas nas figuras em região sombreada, mas não sustentam nenhuma
conclusão.

Como referência temporal para o cálculo da antecedência utiliza-se uma âncora
definida sobre a acurácia de validação --- a época em que ela passa a cair de
forma sustentada. A opção pela acurácia, e não pela perda, decorre de uma
observação empírica relatada no Capítulo~\ref{cap:resultados}: em repositórios
que saturam rapidamente, a perda de validação começa a subir muito cedo, mesmo
em execuções sem memorização, o que a torna uma âncora pouco discriminativa.

A partição dos dados segue a separação em três conjuntos disjuntos. O conjunto
de treino é o único que participa da retropropagação; o de validação é
observado a cada época e nunca entra no gradiente; o de teste é consultado uma
única vez, ao final, na época selecionada.

\section{Síntese}
\label{sec:T-sintese}

Os conceitos aqui reunidos sustentam a proposta em três pontos. A complexidade
LMC oferece uma leitura da distribuição de valores dos pesos, insensível à
posição. A entropia amostral, uni e bidimensional, oferece uma leitura sensível
à organização --- e, no caso bidimensional, preserva a estrutura de linhas e
colunas da matriz de pesos. A formalização do caminho e da ordem de achatamento
torna explícita uma escolha que, de outro modo, permaneceria arbitrária e não
reprodutível.
